\chapter{Noise, SNR, and Random Signals}

Up to this chapter we have treated signals as if they were exact. A sinusoid
was \emph{the} sinusoid; a sampled waveform was a clean list of numbers we
could push through a filter and trust. Reality is messier. Plug a microphone in
and listen to ``silence'' --- you will hear a hiss. Wire an accelerometer to a
microcontroller, hold it perfectly still, and the readings still jitter in the
last few bits. Press two electrodes to your forearm to pick up the tiny
voltages of a muscle and you will find the muscle's signal half-buried under
hum, drift, and crackle. Every real measurement arrives as
\[
  \text{what you wanted} \;+\; \text{what you did not.}
\]
That second term is \emph{noise}, and learning to think about it clearly is not
a sad footnote to signal processing --- it is half the subject. The reason the
Fourier transform, filtering, and averaging are worth so much is precisely that
they let us pull a wanted signal out of an unwanted swamp.

The trouble is that noise is, by its nature, not a fixed waveform. You cannot
write down ``the'' noise the way you write down a sine. Run the same experiment
twice and you get two different hiss patterns. So we need a new kind of object:
a \emph{random signal}, a signal we describe not by its exact values but by its
\emph{statistics} --- how big it tends to be, how fast it wiggles, how its value
now relates to its value a moment ago. This is where the companion probability
book pays off directly. Expectation, variance, and correlation, which there were
abstract operations on random variables, become here concrete, audible,
plottable properties of signals. By the end of the chapter you will be able to
say what white noise \emph{is}, why averaging cleans up a measurement and
exactly how much, why a matched filter is the optimal detector, and what kinds
of noise a wrist-worn IMU and forearm EMG actually fight.

\section{Signal plus noise: the basic picture}
\label{ch20:sig-plus-noise}

The workhorse model for almost everything in this chapter is the simplest one
imaginable. The measured signal $x[n]$ is the true signal $s[n]$ we care about,
plus an additive noise term $w[n]$:
\begin{equation}
  x[n] \;=\; s[n] \;+\; w[n].
  \label{ch20:additive}
\end{equation}
``Additive'' means the noise just adds on top; it does not, say, multiply the
signal or warp time. ``Additive'' is not always exactly right --- quantization
error from Chapter~19 is signal-dependent, and some sensor noise scales with the
reading --- but it is right often enough, and simple enough, that it is the place
to start. Two more assumptions usually ride along: the noise has \emph{zero
mean} (it is just as likely to push the reading up as down), and it is
\emph{independent} of the signal (the muscle twitching does not change the
thermal hiss in the wire).

\begin{figure}[H]
\centering
\begin{tikzpicture}
\begin{axis}[
    width=12cm, height=4.2cm,
    xlabel={sample $n$}, ylabel={amplitude},
    xmin=0, xmax=6.3, ymin=-2.2, ymax=2.2,
    axis lines=left, samples=160, domain=0:6.3,
    title={A clean signal (blue) and the same signal buried in noise (red)},
    every axis plot/.append style={line width=0.8pt},
]
\addplot[specred, very thick] {sin(deg(2*x)) + 0.55*sin(deg(11*x+1)) - 0.45*cos(deg(7*x)) + 0.5*sin(deg(17*x))};
\addplot[sigblue, very thick] {sin(deg(2*x))};
\end{axis}
\end{tikzpicture}
\caption{The same underlying slow tone $s$ shown clean (blue) and
after a noisy channel adds a hash of high-frequency junk (red). The eye still
finds the slow swell; that is a hint that frequency separation will help us.}
\label{ch20:fig-sigplusnoise}
\end{figure}

\begin{intuitionbox}
Why is additive zero-mean noise such a friendly enemy? Because both of its
defining features point to a cure. ``Additive'' means the noise rides
\emph{alongside} the signal rather than mangling it, so anything linear we do to
$x$ --- a filter, a sum, an average --- acts on $s$ and $w$ separately. ``Zero
mean'' means the noise has nowhere to hide on average: pile up enough
independent copies and the positive and negative excursions cancel, while the
signal, repeating itself, reinforces. Almost every noise-fighting trick in this
chapter is one of those two observations cashed in.
\end{intuitionbox}

\subsection{We cannot write down the noise, so we describe its habits}

Here is the conceptual jump. Suppose I record the noise on my EMG amplifier for
one second, twice in a row, with the muscle relaxed. The two recordings are
different number-for-number --- yet they ``look the same'': same typical
loudness, same buzzy texture, same lack of any slow trend. What is stable across
runs is not the values but their \emph{statistics}. A random signal (the proper
name is a \emph{stochastic process}) is best thought of as a whole
\emph{ensemble} of possible waveforms, one of which nature hands you on any given
day. We describe the ensemble, not the single draw.

\begin{connectionbox}
This is exactly the random-variable viewpoint from the probability book, with a
clock attached. There, a random variable $W$ had a distribution; we summarized
it by its mean $\E[W]$ and variance $\Var(W)$. Here we have a random variable
$w[n]$ \emph{at every time index $n$}, an entire indexed family of them. The new
ingredient --- the thing time buys us --- is that we can also ask how
$w[n]$ relates to $w[m]$ at a \emph{different} time. That relationship, the
autocorrelation, is the genuinely new statistic of this chapter, and it is what
the power spectrum will turn out to be.
\end{connectionbox}

\section{Mean, power, and variance of a random signal}
\label{ch20:mean-power}

Two numbers describe ``where a random signal sits'' and ``how much it wiggles.''
The \emph{mean} is the ensemble average at each time,
\begin{equation}
  \mu_x[n] \;=\; \E\bigl[x[n]\bigr],
  \label{ch20:mean}
\end{equation}
where the expectation $\E[\cdot]$ averages over the ensemble --- imagine running
the experiment infinitely many times and averaging the value you get at index
$n$. For our additive model with zero-mean noise, $\mu_x[n] = s[n]$: the mean
\emph{is} the clean signal. That is already a profound statement. It says that if
we could only average over enough independent runs, the noise would melt away and
leave $s$ behind. Section~\ref{ch20:averaging} makes that quantitative.

The second number is \emph{average power}. For a deterministic discrete signal,
power is the mean of the squared values. For a random signal it is the
\emph{expected} squared value,
\begin{equation}
  P_x \;=\; \E\bigl[x[n]^2\bigr].
  \label{ch20:power}
\end{equation}
The names come from circuits: a voltage $v$ across a $1\,\Omega$ resistor
dissipates $v^2$ watts, so $\E[v^2]$ is literally average power. We carry the
word over to every kind of signal. When the mean is zero, power and
\emph{variance} coincide,
\[
  \Var(x[n]) = \E\!\left[(x[n]-\mu_x[n])^2\right] = \E[x[n]^2] = P_x,
\]
and we usually write the noise power as $\sigma_w^2 = \E[w[n]^2]$. Because the
signal and noise are independent and the noise is zero-mean, their powers
\emph{add}:
\begin{equation}
  P_x \;=\; \underbrace{\E[s[n]^2]}_{P_s} \;+\; \underbrace{\E[w[n]^2]}_{\sigma_w^2}.
  \label{ch20:power-add}
\end{equation}
The cross term $2\,\E[s[n]\,w[n]]$ vanishes because $s$ is fixed and $w$ has zero
mean (or, if $s$ is itself random, because they are independent). Powers adding,
not amplitudes, is the single most useful bookkeeping fact about noise.

\begin{whybox}
Why do \emph{powers} add rather than amplitudes? Geometry. Think of the signal
and the noise as two vectors. Independence with zero mean makes them
\emph{orthogonal} on average: their inner product $\E[s\,w]$ is zero. And for
orthogonal vectors the Pythagorean theorem holds --- the squared length of the
sum is the sum of the squared lengths. ``Power'' is squared length, so
$P_x = P_s + \sigma_w^2$ is just Pythagoras in the space of random signals. This
is the same orthogonality that ran through the linear-algebra book; noise simply
gives it a probabilistic flavor.
\end{whybox}

\subsection{Stationarity and ergodicity: when one recording is enough}

There is a practical worry lurking. The mean and power above are
\emph{ensemble} averages --- averages over many parallel universes of the
experiment --- but in the lab you usually have \emph{one} recording. When can a
\emph{time} average over your single record stand in for the ensemble average you
cannot perform?

Two assumptions rescue us. A process is (wide-sense) \emph{stationary} if its
statistics do not drift with time: the mean is a constant $\mu_x$, and the power
and correlations depend only on time \emph{differences}, not on absolute time. A
stationary process is \emph{ergodic} if, in addition, a long enough time average
along one record converges to the ensemble average. For ergodic stationary noise,
\[
  \frac{1}{N}\sum_{n=0}^{N-1} x[n] \;\xrightarrow[N\to\infty]{}\; \mu_x,
  \qquad
  \frac{1}{N}\sum_{n=0}^{N-1} x[n]^2 \;\xrightarrow[N\to\infty]{}\; P_x.
\]

\begin{trapbox}
Ergodicity is an assumption, not a guarantee, and real sensors violate it
constantly. An IMU warming up drifts (its mean moves --- not stationary). EMG
during a contraction is louder than at rest (its power changes with time --- not
stationary). Treating a non-stationary record as if a single long time average
described it is a classic way to compute a confident, wrong number. The safe
move is to work on short windows over which stationarity is plausible --- which is
exactly why the short-time Fourier transform of Chapter~10 exists.
\end{trapbox}

\section{Autocorrelation: how a signal relates to its own past}
\label{ch20:autocorr}

Mean and power tell you the size of a random signal but nothing about its
\emph{texture}. Two noises can have identical power yet feel completely
different: one a smooth low rumble, the other a sharp hiss. The difference is in
how each sample relates to its neighbors. A rumble changes slowly, so $x[n]$ and
$x[n{+}1]$ are nearly equal --- strongly correlated. A hiss jumps around, so
neighbors are nearly unrelated. The statistic that captures this is the
\emph{autocorrelation}: the correlation of the signal with a time-shifted copy
of itself.

For a zero-mean wide-sense-stationary process, define the autocorrelation at lag
$k$ as
\begin{equation}
  R_x[k] \;=\; \E\bigl[\,x[n]\,x[n+k]\,\bigr].
  \label{ch20:Rdef}
\end{equation}
It depends only on the lag $k$, not on $n$ (that is what stationarity buys). A
few properties read straight off the definition and are worth committing to
intuition:

\begin{itemize}
  \item $R_x[0] = \E[x[n]^2] = P_x$. \emph{The autocorrelation at zero lag is the
        power.} A signal is always perfectly correlated with itself.
  \item $R_x[-k] = R_x[k]$: the autocorrelation is symmetric. Looking back and
        looking forward by the same amount give the same correlation.
  \item $\abs{R_x[k]} \le R_x[0]$: the peak is always at lag zero. Nothing
        correlates with the signal better than the signal itself.
  \item How fast $R_x[k]$ \emph{decays} as $k$ grows measures how quickly the
        signal forgets its past. Slow decay $\Rightarrow$ smooth, low-frequency
        signal; fast decay $\Rightarrow$ rough, high-frequency signal.
\end{itemize}

\begin{figure}[H]
\centering
\begin{tikzpicture}
\begin{axis}[
    width=6.0cm, height=4cm, name=leftax,
    xlabel={lag $k$}, ylabel={$R_x[k]$},
    xmin=-8, xmax=8, ymin=-0.4, ymax=1.1,
    axis lines=left, title={White noise}, ytick={0,1},
]
\addplot[ycomb, specred, thick, mark=*, mark size=1.3pt] coordinates {
  (-8,0)(-7,0)(-6,0)(-5,0)(-4,0)(-3,0)(-2,0)(-1,0)(0,1)(1,0)(2,0)(3,0)(4,0)(5,0)(6,0)(7,0)(8,0)};
\end{axis}
\begin{axis}[
    width=6.0cm, height=4cm, at={($(leftax.east)+(1.4cm,0)$)}, anchor=west,
    xlabel={lag $k$}, ylabel={$R_x[k]$},
    xmin=-8, xmax=8, ymin=-0.4, ymax=1.1,
    axis lines=left, title={Smooth (colored) noise}, ytick={0,1},
]
\addplot[ycomb, sigblue, thick, mark=*, mark size=1.3pt] coordinates {
  (-8,0.06)(-7,-0.09)(-6,-0.14)(-5,-0.05)(-4,0.12)(-3,0.30)(-2,0.55)(-1,0.82)(0,1)(1,0.82)(2,0.55)(3,0.30)(4,0.12)(5,-0.05)(6,-0.14)(7,-0.09)(8,0.06)};
\end{axis}
\end{tikzpicture}
\caption{Autocorrelation is the fingerprint of texture. White noise (left) is a
single spike at lag~$0$: each sample is uncorrelated with every other. Smooth
colored noise (right) has a broad hump --- neighboring samples are similar, so
the signal remembers its past for several lags.}
\label{ch20:fig-autocorr}
\end{figure}

\begin{happeningbox}
Autocorrelation is just an inner product with a sliding copy of yourself. Recall
from Chapter~3 that the inner product $\inner{a}{b}$ measures how much two
vectors point the same way. $R_x[k]$ is $\inner{x[n]}{x[n+k]}$ averaged over the
ensemble: it asks ``after I slide the signal by $k$, how much does it still line
up with the original?'' At $k=0$ it lines up perfectly (that is the power). As
the shift grows, a rough signal falls out of step quickly and a smooth one stays
aligned. The whole curve is a self-portrait of the signal's memory.
\end{happeningbox}

\begin{example}[Autocorrelation of a noisy sinusoid]
\label{ch20:ex-noisy-sine}
Let $x[n] = A\cos(\omega_0 n + \phi) + w[n]$, where $\phi$ is a random phase
uniform on $[0,2\pi)$ and $w$ is independent zero-mean white noise with power
$\sigma_w^2$. Using $\E[\cos\alpha\cos\beta]=\tfrac12\cos(\alpha-\beta)$ after
averaging over $\phi$,
\[
  R_x[k] \;=\; \frac{A^2}{2}\cos(\omega_0 k) \;+\; \sigma_w^2\,\delta[k].
\]
The sinusoid contributes a cosine in the lag $k$ that \emph{never dies} --- a
periodic signal remembers its phase forever --- while the noise contributes a
single spike at $k=0$. This is the key to finding a periodic signal in noise: at
nonzero lags the noise has averaged away and only the cosine survives. Pitch
detectors and heart-rate monitors exploit exactly this.
\end{example}

\begin{labbox}
Generate \texttt{x = np.sin(0.3*np.arange(2000)) + np.random.randn(2000)} and
plot \texttt{np.correlate(x, x, 'full')} near the center. Even though the
sinusoid is invisible in the raw trace (the noise is as big as the signal), the
autocorrelation shows clean ripples at the signal's period. You are watching the
noise spike collapse to lag~$0$ while the periodic memory stands proud.
\end{labbox}

\section{Power spectral density: where the power lives in frequency}
\label{ch20:psd}

We now have two ways to describe a random signal's structure: its
autocorrelation $R_x[k]$ in the lag domain, and --- waiting in the wings --- its
spectrum in the frequency domain. The bridge between them is one of the most
important results in all of signal processing, and it is exactly the kind of
time/frequency duality this book keeps returning to.

The \emph{power spectral density} (PSD) $S_x(e^{\jj\omega})$ of a stationary
random signal tells you how much of the signal's power sits at each frequency.
You cannot take an ordinary Fourier transform of a noise waveform --- it never
settles down, so the transform does not converge. But you \emph{can} transform
its autocorrelation, which is a nice decaying deterministic sequence. That is the
content of the Wiener--Khinchin theorem.

\begin{formalbox}
\textbf{Wiener--Khinchin theorem.} For a wide-sense-stationary process, the power
spectral density is the discrete-time Fourier transform of the autocorrelation:
\begin{equation}
  S_x(e^{\jj\omega}) \;=\; \sum_{k=-\infty}^{\infty} R_x[k]\,e^{-\jj\omega k},
  \qquad
  R_x[k] \;=\; \frac{1}{2\pi}\int_{-\pi}^{\pi} S_x(e^{\jj\omega})\,e^{\jj\omega k}\,\dd\omega.
  \label{ch20:wk}
\end{equation}
Setting $k=0$ in the inverse relation gives the total power as the area under the
PSD:
\[
  P_x = R_x[0] = \frac{1}{2\pi}\int_{-\pi}^{\pi} S_x(e^{\jj\omega})\,\dd\omega.
\]
\end{formalbox}

\begin{whybox}
Why should the spectrum of power be the transform of the autocorrelation, of all
things? Think about what each frequency component contributes. A signal that
oscillates strongly at frequency $\omega_0$ has an autocorrelation that itself
oscillates at $\omega_0$ (Example~\ref{ch20:ex-noisy-sine}). The Fourier
transform's whole job is to detect ``how much oscillation at each frequency'' is
present in a sequence. Apply it to $R_x[k]$ and it reports how much oscillation
the signal's memory contains at each $\omega$ --- which is precisely how the power
is distributed across frequency. The transform of the self-similarity is the
spectrum of the energy. Wiener and Khinchin made this airtight, but the picture
is that simple.
\end{whybox}

\begin{connectionbox}
The probability book defined the variance of a sum of independent contributions
as the sum of variances. The PSD is that statement made spectral: it splits the
total variance (power) into a density across frequency, $\Var = \int S_x\,\dd f$.
Each frequency band is an independent ``bucket'' of variance, and they add. When
later you read that a filter ``shapes the noise spectrum,'' it means it
reweights those buckets before they are added back up.
\end{connectionbox}

\section{White noise and colored noise}
\label{ch20:white-colored}

We can now define the two most important species of noise by the shape of their
spectra --- and the names are not metaphors but a direct analogy with light.

\textbf{White noise} has a \emph{flat} PSD: equal power at every frequency, just
as white light contains every color of the visible spectrum equally. Flat in
frequency means, by Wiener--Khinchin, an autocorrelation that is a single spike,
\[
  S_w(e^{\jj\omega}) = \sigma_w^2 \;\;(\text{constant})
  \quad\Longleftrightarrow\quad
  R_w[k] = \sigma_w^2\,\delta[k].
\]
A spike at lag $0$ and zero elsewhere says: every sample is uncorrelated with
every other sample. White noise has no memory at all. It is the most ``random''
a signal can be, and it is the standard idealized model for thermal (Johnson)
noise in resistors and shot noise in electronics.

\begin{trapbox}
``White'' says nothing about the \emph{distribution} of the values --- only that
samples are uncorrelated. White Gaussian noise (samples drawn from a bell curve)
is the famous case, but you can have white noise whose samples are $\pm 1$ coin
flips. Conversely, ``Gaussian'' says nothing about correlation; you can have
Gaussian noise that is heavily colored. Whiteness is about \emph{the spectrum};
Gaussianity is about \emph{the histogram}. Keep the two axes separate.
\end{trapbox}

\textbf{Colored noise} has a PSD that is not flat: its power is concentrated in
some bands more than others, so successive samples \emph{are} correlated.
``Coloring'' white noise just means filtering it (Section~\ref{ch20:filtered}).
Common colors, named again after light:
\begin{itemize}
  \item \emph{Pink noise} ($S(f)\propto 1/f$): more power at low frequencies,
        falling $3$~dB per octave. Ubiquitous in nature --- it is the spectrum of
        the rumble of traffic, the flicker noise in electronics, and slow sensor
        drift.
  \item \emph{Brown / red noise} ($S(f)\propto 1/f^2$): even more low-frequency
        dominance; the spectrum of a random walk. This is what integrating white
        noise produces --- and integrating a gyroscope's white rate noise is
        \emph{exactly} how orientation estimates drift.
  \item \emph{Blue / violet noise}: power rising with frequency, the spectrum you
        get from \emph{differentiating} a signal. Numerical differentiation of a
        noisy sensor amplifies high-frequency noise for this reason.
\end{itemize}

\begin{figure}[H]
\centering
\begin{tikzpicture}
\begin{axis}[
    width=12cm, height=5cm,
    xlabel={frequency $f$ (normalized)}, ylabel={PSD $S(f)$ (relative)},
    xmin=0.02, xmax=1, ymin=0.05, ymax=30,
    xmode=log, ymode=log, axis lines=left, samples=80, domain=0.02:1,
    legend pos=south west, legend cell align=left,
    title={Spectra of white, pink, and brown noise},
    every axis plot/.append style={line width=1pt},
]
\addplot[gray, very thick] {1}; \addlegendentry{white ($f^{0}$)}
\addplot[sigblue, very thick] {0.6/x}; \addlegendentry{pink ($1/f$)}
\addplot[quotientorange, very thick] {0.12/(x^2)}; \addlegendentry{brown ($1/f^2$)}
\end{axis}
\end{tikzpicture}
\caption{Noise ``colors'' on log-log axes. White is flat. Pink falls as $1/f$
(a straight line of slope $-1$ here). Brown falls as $1/f^2$ (slope $-2$),
piling almost all of its power into the lowest frequencies --- which is why
random-walk drift dominates the long-term behavior of integrated sensors.}
\label{ch20:fig-colors}
\end{figure}

\begin{historybox}
The ``white'' metaphor traces to the recognition, in the 1920s, that thermal
agitation in a conductor produces a voltage with an essentially flat spectrum ---
John B. Johnson measured it at Bell Labs in 1928 and Harry Nyquist explained it
that same year from thermodynamics, giving the formula $S_v = 4 k_B T R$ for the
noise power spectral density of a resistor $R$ at temperature $T$. Norbert Wiener
(1930) and Aleksandr Khinchin (1934) independently established the
autocorrelation--spectrum link that lets us even speak of the ``spectrum'' of a
signal that never repeats. Together these results turned noise from a nuisance
into a quantity you can compute, predict, and design against.
\end{historybox}

\subsection{Coloring noise is just filtering it}
\label{ch20:filtered}

Pass white noise with power $\sigma_w^2$ through an LTI filter with frequency
response $H(e^{\jj\omega})$. What comes out? Because filtering multiplies each
frequency component by $H$, and power is amplitude squared, the output PSD is the
input PSD scaled by $\abs{H}^2$:
\begin{equation}
  S_y(e^{\jj\omega}) \;=\; \abs{H(e^{\jj\omega})}^2\, S_w(e^{\jj\omega})
  \;=\; \abs{H(e^{\jj\omega})}^2\,\sigma_w^2.
  \label{ch20:filtered-psd}
\end{equation}
This single formula is the engine behind every ``colored'' noise: choose the
filter shape and you choose the color. A one-pole low-pass filter applied to
white noise yields pink-ish noise; an integrator ($H \approx 1/(1-z^{-1})$)
yields brown noise; a differencer yields blue noise. It also tells you exactly
how a filter changes the noise power: integrate $\abs{H}^2\sigma_w^2$ over
frequency. A filter that passes only a tenth of the band lets through only a
tenth of white-noise power --- the basis of using a band-pass filter to clean up a
narrowband measurement.

\begin{example}[Noise out of a moving-average filter]
\label{ch20:ex-ma-noise}
The simplest smoother, the $L$-point moving average $y[n]=\frac1L\sum_{i=0}^{L-1}x[n-i]$,
fed white noise of power $\sigma_w^2$, produces output power
\[
  \sigma_y^2 = \frac{1}{L^2}\sum_{i=0}^{L-1}\sigma_w^2 = \frac{\sigma_w^2}{L}.
\]
The noise power drops by a factor of $L$ (the amplitude by $\sqrt L$). The
filter has thrown away the high-frequency part of the white spectrum, keeping
only a low-pass slice. We are about to see this same $1/L$ law from the averaging
side --- they are the same fact.
\end{example}

\section{Signal-to-noise ratio (SNR)}
\label{ch20:snr}

We need one number to say ``how clean is this measurement?'' That number is the
\emph{signal-to-noise ratio}, the ratio of signal power to noise power:
\begin{equation}
  \mathrm{SNR} \;=\; \frac{P_s}{\sigma_w^2}
  \;=\; \frac{\E[s[n]^2]}{\E[w[n]^2]}.
  \label{ch20:snr-def}
\end{equation}
Because powers can span enormous ranges, SNR is almost always quoted in
\emph{decibels}:
\begin{equation}
  \mathrm{SNR_{dB}} \;=\; 10\log_{10}\frac{P_s}{\sigma_w^2}
  \;=\; 20\log_{10}\frac{A_s}{A_w},
  \label{ch20:snr-db}
\end{equation}
where $A_s, A_w$ are RMS amplitudes (the factor changes from $10$ to $20$ because
power is amplitude squared, and $\log$ of a square is twice the log). A few
anchor points worth memorizing: $\mathrm{SNR}=1$ means $0$~dB (signal and noise
equally loud); every factor of $10$ in power is $+10$~dB; every doubling of
\emph{amplitude} is about $+6$~dB; a factor of $2$ in power is about $+3$~dB.

\begin{intuitionbox}
Decibels feel fussy until you notice they turn the two operations you care about
into addition. Cascading two stages \emph{multiplies} their power gains, which
\emph{adds} their dB values. And the human ear and eye both respond roughly
logarithmically, so dB also matches perception: a steady $+10$~dB feels like ``a
notch louder'' whether you start from a whisper or a shout. SNR in dB is the
natural currency for talking about how much headroom a measurement has.
\end{intuitionbox}

\begin{example}[Reading SNR off a recording]
\label{ch20:ex-snr-read}
Suppose a relaxed-muscle EMG baseline has RMS $0.01$~mV and a contraction
produces RMS $0.20$~mV. The signal (the part above baseline) sits at roughly
$0.20$~mV against $0.01$~mV of noise, so
\[
  \mathrm{SNR_{dB}} = 20\log_{10}\frac{0.20}{0.01} = 20\log_{10}(20) \approx 26\ \text{dB}.
\]
That is a comfortably detectable contraction. At $0$~dB (noise as loud as
signal) you would struggle; below that, you need the averaging and filtering of
the next sections, or a better electrode.
\end{example}

\begin{trapbox}
SNR is meaningless until you say \emph{over what bandwidth} and \emph{for what
signal}. White noise has power spread over the whole band, so the noise that
actually hurts you is only the part inside your signal's band. Quoting an SNR
without naming the band (or the filter that defines it) is like quoting a speed
without a direction. The right SNR for detection is computed \emph{after} the
processing that isolates the signal --- which is precisely why filtering and
matched filtering raise the SNR you get to use.
\end{trapbox}

\section{Averaging to beat down noise}
\label{ch20:averaging}

Here is the most-used noise weapon in all of measurement, and it follows
directly from ``powers add'' plus ``the mean is the clean signal.'' Suppose you
can record the \emph{same} signal $s[n]$ many times --- $M$ repetitions, each with
independent noise:
\[
  x_m[n] = s[n] + w_m[n], \qquad m = 1,\dots,M.
\]
Average the recordings sample by sample:
\[
  \bar x[n] = \frac{1}{M}\sum_{m=1}^{M} x_m[n]
            = s[n] + \frac{1}{M}\sum_{m=1}^{M} w_m[n].
\]
The signal term is untouched --- it is identical in every run, so averaging
$M$ copies of $s[n]$ just gives $s[n]$ back. The noise term is an average of $M$
\emph{independent} zero-mean noises. Its power is
\[
  \Var\!\left(\frac{1}{M}\sum_m w_m[n]\right)
  = \frac{1}{M^2}\sum_m \Var(w_m[n])
  = \frac{1}{M^2}\cdot M\sigma_w^2
  = \frac{\sigma_w^2}{M}.
\]
The noise power falls by a factor of $M$; the noise \emph{amplitude} (its RMS)
falls by $\sqrt{M}$. And so:
\begin{equation}
  \mathrm{SNR}_{\bar x} = M\cdot \mathrm{SNR}_x,
  \qquad
  \mathrm{SNR_{dB}}\ \text{improves by}\ 10\log_{10} M.
  \label{ch20:avg-gain}
\end{equation}

\begin{figure}[H]
\centering
\begin{tikzpicture}
\begin{axis}[
    width=11cm, height=5cm,
    xlabel={number of averages $M$}, ylabel={SNR improvement (dB)},
    xmin=1, xmax=1000, ymin=0, ymax=32,
    xmode=log, log basis x=10, axis lines=left, samples=60, domain=1:1000,
    title={Averaging $M$ records buys $10\log_{10}M$ dB --- diminishing returns},
    every axis plot/.append style={line width=1.1pt},
    ymajorgrids=true,
]
\addplot[sigblue, very thick] {10*ln(x)/ln(10)};
\node[anchor=west, font=\small] at (axis cs:1.2,4){$M{=}1$: $0$ dB};
\node[anchor=west, font=\small] at (axis cs:9,13){$M{=}10$: $+10$ dB};
\node[anchor=east, font=\small] at (axis cs:900,27){$M{=}1000$: $+30$ dB};
\end{axis}
\end{tikzpicture}
\caption{The $\sqrt{M}$ law has a sting in its tail. The first ten averages buy
a full $10$~dB, but to gain the \emph{next} $10$~dB you need ten times as many
averages again. Quadrupling your averaging time only halves the noise. Beyond a
point, a quieter sensor beats more averaging.}
\label{ch20:fig-avg}
\end{figure}

\begin{happeningbox}
Why does the signal survive averaging while the noise dies? Because they are
\emph{coherent} and \emph{incoherent}, respectively. The signal is identical in
every run, so its copies add ``in phase'' --- $M$ copies of amplitude $A$ give
amplitude $MA$, i.e.\ power grows as $M^2$. The noise is independent run to run,
so its copies add ``out of phase'' like a random walk --- $M$ steps give a typical
distance $\sqrt{M}$, i.e.\ power grows only as $M$. Divide both by $M$ (the
averaging) and the signal sits at constant amplitude while the noise shrinks as
$1/\sqrt{M}$. Coherent gain beats incoherent gain; that gap \emph{is} the SNR
improvement.
\end{happeningbox}

\begin{trapbox}
Averaging only works when the signal is genuinely repeatable and the runs can be
\emph{aligned}. If the signal jitters in time between runs (different trigger
delays), naive averaging smears and attenuates it --- you average away the very
thing you wanted. This is why evoked-potential and radar systems trigger so
carefully, and why gesture pipelines time-align repetitions before averaging
templates. Also: the $1/M$ law assumes \emph{independent} noise. Correlated noise
(a shared 50/60~Hz hum on every run) does not average down at all --- it is
coherent too. Averaging fights random noise, not interference.
\end{trapbox}

\subsection{Two flavors: ensemble averaging vs.\ smoothing}

The averaging above was \emph{across repetitions} (ensemble averaging). When you
cannot repeat the experiment, you can instead average \emph{across time} within
one record --- a moving average, which is just a low-pass filter. The
$\sigma_w^2/L$ result of Example~\ref{ch20:ex-ma-noise} is the same $1/M$ law in
disguise, with $M=L$ neighboring samples playing the role of repetitions. The
price is different though: ensemble averaging keeps the signal intact, while time
smoothing also low-passes the signal, blurring any fast features it had. There is
no free lunch --- you trade noise for bandwidth.

\begin{example}[How long to average?]
\label{ch20:ex-howlong}
You measure a static accelerometer offset buried in white noise of RMS
$5\,\text{mg}$ and you want the estimate good to $0.5\,\text{mg}$ RMS. You need
the noise to fall by a factor of $10$ in amplitude, hence by $100$ in power, so
$M = 100$ samples. At $100$~Hz output rate that is one second of stillness.
Want $0.05\,\text{mg}$? That is $\times 100$ in amplitude, $M=10{,}000$ samples,
$100$ seconds --- and by then the bias has probably drifted (non-stationarity),
so more averaging stops helping. The math tells you when to stop chasing it.
\end{example}

\section{Matched filtering: the optimal detector}
\label{ch20:matched}

Averaging assumed you could repeat the experiment. Often you cannot --- you get
one noisy record and must decide \emph{whether and where} a known signal shape
appears in it: a radar pulse echoing back, a sync word in a data packet, a
particular flick of the wrist in an IMU stream. The optimal tool is the
\emph{matched filter}, and its logic is pure geometry.

Suppose the known signal shape (the ``template'') is $s[n]$, of length $L$, and
your record is $x[n]=s[n-n_0]+w[n]$ with white noise --- the template is hiding at
unknown position $n_0$. To test position $m$, slide the template there and take
the inner product with the data:
\begin{equation}
  y[m] \;=\; \sum_{n} s[n-m]\,x[n] \;=\; (x \conv \tilde s)[m],
  \label{ch20:matched-corr}
\end{equation}
which is a correlation --- equivalently a convolution with the
\emph{time-reversed} template $\tilde s[n]=s[-n]$ (Figure~\ref{ch20:fig-matched-block}).
The output $y[m]$ peaks sharply when the template lines up with the buried copy.

\begin{whybox}
Why is correlating with the template the \emph{best} possible thing to do? Cast
it as vectors. The data over a window is the signal vector $\vec s$ plus a noise
vector $\vec w$. A linear detector is an inner product $\inner{\vec h}{\vec x}$
with some weight vector $\vec h$. The signal part of the output is
$\inner{\vec h}{\vec s}$ and the noise part has power $\sigma_w^2\,\norm{\vec h}^2$
(white noise). We want to maximize output SNR,
$\inner{\vec h}{\vec s}^2 \big/ (\sigma_w^2\norm{\vec h}^2)$. By the
Cauchy--Schwarz inequality, $\inner{\vec h}{\vec s}^2 \le \norm{\vec h}^2\norm{\vec s}^2$,
with equality \emph{only when $\vec h$ is parallel to $\vec s$}. So the optimal
weights \emph{are} the signal itself: $\vec h = \vec s$. The matched filter wins
because it points the detector exactly along the signal direction, collecting
all of the signal's energy while seeing noise from only that one direction.
\end{whybox}

\begin{figure}[H]
\centering
\begin{tikzpicture}
\begin{axis}[
    width=12cm, height=4.6cm,
    xlabel={lag $m$}, ylabel={output $y[m]$},
    xmin=0, xmax=12, ymin=-1.2, ymax=3.4,
    axis lines=left, samples=130, domain=0:12,
    title={The matched filter turns a buried pulse into a clean peak},
    every axis plot/.append style={line width=1pt},
]
\addplot[gray, thick] {0.45*sin(deg(6*x)) + 0.4*cos(deg(9*x+1)) - 0.35*sin(deg(13*x))};
\addplot[specred, very thick] {0.45*sin(deg(6*x)) + 0.4*cos(deg(9*x+1)) - 0.35*sin(deg(13*x)) + 3*exp(-((x-7)*1.1)^2)};
\end{axis}
\end{tikzpicture}
\caption{Output of a matched filter sliding along a noisy record (gray: the
noise-only baseline). Where the template aligns with the hidden copy (here near
$m=7$) the correlation collects all the signal energy into one sharp peak (red)
that stands far above the noise floor. Detection becomes ``is there a peak above
threshold?''}
\label{ch20:fig-matched}
\end{figure}

\begin{happeningbox}
The peak height of a matched filter is the signal's \emph{energy}
$\sum_n s[n]^2$, while the noise at the output has power $\sigma_w^2\sum_n s[n]^2$.
So the output SNR is the total signal \emph{energy} divided by the noise PSD ---
it depends only on how much energy the pulse carries, \emph{not} on its shape.
A long quiet chirp and a short loud spike with the same energy detect equally
well. That is the deep reason radar uses long coded pulses: spread the energy out
in time to stay within power limits, then let the matched filter squeeze it back
into a sharp peak. The technique is called \emph{pulse compression}.
\end{happeningbox}

\begin{figure}[H]
\centering
\begin{tikzpicture}[
    box/.style={draw, rounded corners, minimum height=8mm, minimum width=20mm,
                font=\small, fill=blue!4},
    >=Stealth, node distance=14mm,
]
\node (in) {$x[n]=s[n-n_0]+w[n]$};
\node[box, right=16mm of in] (filt) {$h[n]=s[-n]$};
\node[box, right of=filt, node distance=30mm] (peak) {peak detect};
\node[right=14mm of peak] (out) {decision};
\draw[->] (in) -- (filt);
\draw[->] (filt) -- node[above, font=\scriptsize]{$y[m]$} (peak);
\draw[->] (peak) -- node[above, font=\scriptsize]{$>\!\theta?$} (out);
\end{tikzpicture}
\caption{The matched-filter detector as a block diagram. Convolving the noisy
record with the time-reversed template $h[n]=s[-n]$ correlates the data against
the known shape; the output $y[m]$ is then compared to a threshold $\theta$.
Detection reduces to ``did the correlation peak clear the threshold, and where?''}
\label{ch20:fig-matched-block}
\end{figure}

\begin{connectionbox}
The matched filter is the convolution of Chapter~13 and the inner-product
geometry of Chapter~3 meeting probability. Correlation \emph{is} convolution with
a reversed template; the Cauchy--Schwarz optimality \emph{is} the orthogonality
geometry of inner-product spaces; and ``maximize SNR'' \emph{is} the
expectation/variance machinery of the probability book. Three threads of this
series tie together in one estimator.
\end{connectionbox}

\section{Sensor noise models for the wrist tracker}
\label{ch20:sensors}

Everything above becomes concrete the moment you wire up a real sensor. Here is a
plain inventory of what an IMU and an EMG front end throw at you, in the
language of this chapter.

\subsection{IMU noise: white rate noise and the random walk it integrates to}

A MEMS gyroscope and accelerometer each report a true value plus several noise
terms. The two that matter most:
\begin{itemize}
  \item \textbf{Noise density (white noise).} The instantaneous output has a
        roughly white, zero-mean fluctuation. Datasheets quote it as a
        \emph{density}: e.g.\ a gyro at $0.01\,^{\circ}/\text{s}/\sqrt{\text{Hz}}$.
        The $\sqrt{\text{Hz}}$ is Wiener--Khinchin in disguise: white power is
        spread per unit bandwidth, so the RMS you actually see is the density
        times $\sqrt{\text{bandwidth}}$. Halve your bandwidth (filter harder) and
        the white-noise RMS drops by $\sqrt 2$.
  \item \textbf{Bias instability and random walk (colored noise).} The zero point
        wanders slowly. Integrating the gyro's white rate noise to get angle
        produces \emph{angle random walk} --- brown noise, $S(f)\propto 1/f^2$ ---
        whose error grows as $\sqrt{t}$. A slowly drifting bias on top adds an
        even-lower-frequency wander. This is the unavoidable reason a gyro-only
        orientation estimate drifts and must be corrected by the accelerometer
        and magnetometer (the sensor fusion of Chapter~36).
\end{itemize}

\begin{buildbox}
The standard tool for reading these terms off real data is the \emph{Allan
deviation}: log the IMU at rest for a few hours, then plot the deviation of the
running average versus averaging time $\tau$ on log-log axes. A slope of $-1/2$
at short $\tau$ is the white noise (more averaging helps, the $1/\sqrt M$ law);
the \emph{minimum} of the curve is the bias instability (the floor where more
averaging stops helping); a $+1/2$ slope at long $\tau$ is the random walk
(drift dominating). The $\tau$ at the minimum tells you the longest you should
ever average before recalibrating. A short Python sketch:
\begin{lstlisting}[language=Python]
import numpy as np
def allan_dev(x, fs, taus):
    # x: stationary sensor record; fs: sample rate (Hz)
    out = []
    for tau in taus:
        m = int(tau * fs)                       # samples per averaging bin
        k = len(x) // m
        means = x[:k*m].reshape(k, m).mean(1)   # bin averages
        d = np.diff(means)
        out.append(np.sqrt(0.5 * np.mean(d**2)))  # Allan deviation
    return np.array(out)
\end{lstlisting}
Sample at the highest rate the sensor supports, then low-pass and decimate in
software: averaging in the digital domain is free SNR, whereas noise picked up
before the ADC is gone for good.
\end{buildbox}

\begin{figure}[H]
\centering
\begin{tikzpicture}
\begin{axis}[
    width=11cm, height=5cm,
    xlabel={averaging time $\tau$ (s)}, ylabel={Allan deviation},
    xmin=0.01, xmax=1000, ymin=0.003, ymax=0.3,
    xmode=log, ymode=log, axis lines=left, samples=80, domain=0.01:1000,
    title={Allan deviation of a MEMS gyro: noise, floor, drift},
    every axis plot/.append style={line width=1.2pt},
]
\addplot[sigblue, very thick] {sqrt( (0.02^2)/x + 0.01^2 + (0.0006^2)*x )};
\node[anchor=south west, font=\small, sigblue] at (axis cs:0.012,0.13){slope $-\tfrac12$: white noise};
\node[anchor=south, font=\small, sigblue] at (axis cs:3,0.012){bias floor};
\node[anchor=south east, font=\small, sigblue] at (axis cs:700,0.06){slope $+\tfrac12$: random walk};
\end{axis}
\end{tikzpicture}
\caption{The classic Allan-deviation ``bathtub.'' Short averaging is dominated by
white noise that averages down; a flat minimum marks the bias-instability floor;
long averaging is dominated by random-walk drift that averages \emph{up}. The
dip locates the sweet-spot averaging time and the irreducible noise floor of the
part.}
\label{ch20:fig-allan}
\end{figure}

\subsection{EMG noise: where the real signal lives, and what fights it}

Surface EMG --- the muscle voltages we will use to sense the hand --- is a far
harsher noise problem than the IMU, because the wanted signal is microvolts to
low millivolts and the contaminants are large and structured
(Figure~\ref{ch20:fig-emg}):
\begin{itemize}
  \item \textbf{Powerline interference} at $50/60$~Hz and its harmonics, coupled
        capacitively from mains wiring. This is \emph{not} random noise --- it is a
        coherent sinusoid, so averaging does not touch it. You fight it with a
        notch filter and good differential electrodes.
  \item \textbf{Baseline wander / motion artifact}: low-frequency ($<10$~Hz)
        drift from electrode-skin impedance changes as you move --- colored,
        low-frequency, large. A high-pass filter around $20$~Hz removes most of
        it (at the cost of the lowest EMG content).
  \item \textbf{Thermal and amplifier noise}: broadband, roughly white, set by
        the electrode impedance and the front-end op-amp. This \emph{does}
        average down and sets the ultimate floor.
\end{itemize}
The useful EMG energy sits roughly in the $20$--$450$~Hz band, so the standard
front end is a band-pass ($\approx 20$--$450$~Hz) plus a $50/60$~Hz notch ---
literally Equation~\eqref{ch20:filtered-psd} put to work, shaping
$\abs{H}^2$ to pass the signal band and crush the interference bands.

\begin{figure}[H]
\centering
\begin{tikzpicture}
\begin{axis}[
    width=12cm, height=4.4cm,
    xlabel={time (s)}, ylabel={voltage (mV)},
    xmin=0, xmax=2, ymin=-0.7, ymax=0.7,
    axis lines=left, samples=200, domain=0:2,
    title={A raw surface-EMG trace: a burst hiding under hum and drift},
    every axis plot/.append style={line width=0.6pt},
]
% baseline wander + 50 Hz hum + a contraction burst in the middle
\addplot[sigblue] {0.18*sin(deg(2*pi*0.7*x)) + 0.10*sin(deg(2*pi*50*x))
   + (0.35*exp(-((x-1.0)*2.2)^2))*sin(deg(2*pi*120*x))};
\end{axis}
\end{tikzpicture}
\caption{One second into the record the muscle fires (the fast burst near
$t=1$~s), but the raw trace is dominated by slow baseline wander and steady
$50$~Hz hum. A high-pass kills the wander, a notch kills the hum, and a band-pass
keeps the burst --- the front end of Equation~\eqref{ch20:filtered-psd} at work.}
\label{ch20:fig-emg}
\end{figure}

\begin{goalbox}
For the wrist controller, this chapter sets the rules of the game. The IMU gives
us clean, high-SNR \emph{coarse} motion --- a flick, a tilt, a tap --- because
those gestures live well above the noise floor and we can low-pass away the white
jitter cheaply. But keep in view the limit stated back in the preface: an IMU
on the wrist cannot recover individual finger joint angles. Finger intent shows
up in the forearm \emph{muscles}, and reading it means EMG (or force myography),
with all the powerline, motion-artifact, and microvolt-floor problems above. So
the noise budget drives the whole architecture: use the IMU for robust discrete
gestures where SNR is generous, and reserve EMG --- with aggressive band-pass +
notch filtering and short-window RMS features --- for the finger-level
distinctions, accepting that even then we get coarse muscle-activation classes,
not continuous joint angles. Continuous finger pose needs sensors \emph{on} the
fingers (instrumented rings/gloves). Designing the gesture set to live where the
SNR is good is the single most important decision in the build.
\end{goalbox}

\begin{applicationbox}
The same toolkit recurs across fields under different names. In audio, dithering
deliberately adds white noise to decorrelate quantization error (Chapter~19); in
astronomy, ``stacking'' frames is ensemble averaging that pulls faint galaxies
out of sensor noise; in MRI and EEG, signal averaging over repeated stimuli
recovers evoked responses at SNRs far below $0$~dB; in communications, the
matched filter is the heart of every receiver. Learn the four moves of this
chapter --- model the noise, measure its spectrum, average what is repeatable,
filter/match what is not --- and you can clean up a signal in any domain.
\end{applicationbox}

\section{What to carry forward}
\label{ch20:summary}

\begin{recoverybox}
The compressed version of this chapter:
\begin{itemize}
  \item Real measurements are $x = s + w$. We describe $w$ by statistics, not
        values: a random signal is an \emph{ensemble}.
  \item Mean is the clean signal; \emph{power} (= variance for zero mean) is the
        size of the wiggle. Independent signal and noise have \emph{powers that
        add} --- Pythagoras for random signals.
  \item \emph{Autocorrelation} $R_x[k]$ measures memory; its Fourier transform is
        the \emph{power spectral density} (Wiener--Khinchin). White noise = flat
        PSD = spike autocorrelation; colored noise = shaped PSD = filtered white.
  \item \emph{SNR} $=P_s/\sigma_w^2$, quoted in dB; always say over what band.
  \item \emph{Averaging} $M$ repeatable records improves SNR by $10\log_{10}M$ dB
        ($1/\sqrt M$ in amplitude) --- but only for independent noise and aligned
        signals.
  \item The \emph{matched filter} (correlate with the known template) maximizes
        output SNR; its gain is the signal energy over the noise PSD, by
        Cauchy--Schwarz.
  \item IMU noise is mostly white rate noise plus random-walk drift (Allan
        deviation reveals the split); EMG fights powerline hum, motion artifact,
        and a microvolt floor. Know the limits: IMU $\to$ coarse gestures only;
        finger intent needs EMG/FMG; continuous finger pose needs finger-mounted
        sensors.
\end{itemize}
\end{recoverybox}

\section{Exercises}
\label{ch20:exercises}

\begin{exercise}
A signal has RMS amplitude $2.0$ and is corrupted by independent zero-mean noise
of RMS $0.5$. (a)~What is the SNR as a plain ratio? (b)~In dB? (c)~If you average
$16$ independent recordings, what is the new noise RMS and the new SNR in dB?
\end{exercise}

\begin{exercise}
Prove that for independent, zero-mean noise added to a deterministic signal, the
cross term $\E[s[n]\,w[n]]$ in the power $\E[(s+w)^2]$ is zero, and hence
$P_x = P_s + \sigma_w^2$. State exactly which assumption each step uses.
\end{exercise}

\begin{exercise}
White noise of power $\sigma_w^2=4$ is passed through a filter with
$\abs{H(e^{\jj\omega})}^2 = 2$ for $\abs{\omega}<\pi/4$ and $0$ elsewhere. Using
the area-under-the-PSD formula, compute the output power. (Hint: the input PSD is
flat at $\sigma_w^2$; integrate $\abs{H}^2\sigma_w^2$ over $[-\pi,\pi]$ and divide
by $2\pi$.)
\end{exercise}

\begin{exercise}
Show that the autocorrelation of white noise, $R_w[k]=\sigma_w^2\delta[k]$, has a
\emph{flat} PSD by plugging it into the Wiener--Khinchin sum
\eqref{ch20:wk}. Then do the reverse: given a flat PSD $S(e^{\jj\omega})=c$,
recover $R_w[k]$ via the inverse transform.
\end{exercise}

\begin{exercise}
A random walk is $y[n]=y[n-1]+w[n]$ with white $w$. (a)~Show $y$ is the output of
an integrator $H(z)=1/(1-z^{-1})$ driven by white noise. (b)~Argue from
$\abs{H(e^{\jj\omega})}^2 = 1/\abs{1-e^{-\jj\omega}}^2$ that its PSD behaves like
$1/\omega^2$ near $\omega=0$ --- i.e.\ it is brown noise. (c)~Connect this to why
integrating gyro noise gives drift that grows as $\sqrt t$.
\end{exercise}

\begin{exercise}
You want to detect a known length-$64$ template buried in white noise. (a)~Write
the matched-filter output as a correlation. (b)~If the template has unit energy
and the noise PSD is $\sigma_w^2$, what is the output peak SNR? (c)~If you halve
the template's amplitude everywhere, by how many dB does the peak SNR drop?
\end{exercise}

\begin{exercise}
Explain, using the coherent/incoherent picture, why averaging $M$ recordings does
\emph{not} reduce $50$~Hz powerline interference that is present (in phase) on
every recording, even though it reduces the thermal noise by $\sqrt M$. What
operation \emph{would} remove the powerline component, and why?
\end{exercise}

\begin{exercise}
A gyro datasheet lists noise density $0.014\,^{\circ}/\text{s}/\sqrt{\text{Hz}}$.
(a)~If you run it with a $50$~Hz output bandwidth, estimate the white-noise RMS
in $^{\circ}/\text{s}$. (b)~You apply a moving average of $L=10$ samples at
$200$~Hz sampling. By what factor does the white-noise RMS drop, and what new
effective bandwidth does that correspond to? (c)~Name one noise term that this
averaging will \emph{not} help with.
\end{exercise}

\begin{exercise}
On an Allan-deviation log-log plot, identify which physical noise term produces
each of: a $-1/2$ slope at small $\tau$, a flat minimum, and a $+1/2$ slope at
large $\tau$. Explain in one sentence each why averaging helps, plateaus, then
hurts.
\end{exercise}

\begin{exercise}
(Design.) You are choosing a gesture set for the wrist controller. Given that the
IMU delivers high SNR for coarse motion but cannot resolve finger joints, and EMG
delivers only coarse muscle-activation classes against heavy interference,
propose a $4$-gesture vocabulary and justify each choice by where its SNR is
favorable. State explicitly one gesture you would \emph{not} attempt with this
hardware and why.
\end{exercise}
